% Redrawn 2026-10-15; main labels 8 pt, notes 7.5 pt; actual fonts embedded in PDF.
% export-reference-width-mm: 169
\begin{tikzpicture}[font=\figurefont\fontsize{8}{10}\selectfont,
  st/.style={line width=.6pt,rounded corners=1.8pt,draw=FigureStroke,fill=FigureBlueFill,minimum width=12mm,
    minimum height=9mm,align=center},
  arr/.style={-{Latex[length=1.8mm]},draw=FigureStroke,line width=.6pt}]
  \node[anchor=west,font=\figurefont\mdseries] at (-.6,2.2) {(a) 连续选择才可达};
  \foreach \x/\name/\txt in {0/s0/起点,1.8/s1/门前,3.6/s2/门后,5.4/s3/奖励}{
    \node[st] (\name) at (\x,.9) {\txt};
  }
  \draw[arr] (s0) -- (s1) node[midway,above] {正确};
  \draw[arr] (s1) -- (s2) node[midway,above] {单向};
  \draw[arr] (s2) -- (s3) node[midway,above] {正确};
  \draw[arr,dashed,BookAmber] (s2.south) -- (3.6,-.4)
    -- (0,-.4) -- (s0.south);
  \node[fill=white,text=BookInk] at (1.8,-.4) {错误则重置};
  \node[align=center] at (2.7,-1.3) {每步独立随机选择\\到达率可能很低};
  \begin{scope}[shift={(9.1,0)}]
    \node[anchor=west,font=\figurefont\mdseries] at (-.6,2.2) {(b) 不可控噪声不是新知识};
    \node[st] (pos) at (0,.5) {当前位置};
    \node[st,fill=FigureBlueFill] (door) at (4.3,1.3) {门口\\任务进展};
    \node[st,fill=FigureRedFill] (tv) at (4.3,-.5) {随机显示器\\$\xi=\pm1$};
    \draw[arr] (pos) -- (door) node[midway,above,sloped] {长程行动};
    \draw[arr] (pos) -- (tv) node[midway,below,sloped] {误差奖励};
    \node[align=center] at (2,-1.5) {最优预测为0\\剩余平方误差仍为1};
  \end{scope}
\end{tikzpicture}
